\clearpage\SourceLeaf{207}
\begin{center}
[ÜBER DIE ERSTEN GRÜNDE DER GEOMETRIE.]
\end{center}

\begin{center}
\textsc{Gauss} an \textsc{Bessel}. Göttingen, 27. Januar 1829.
\end{center}

.\ .\ .\ .\ .\ .\ .\ Auch über ein anderes Thema, das mich schon seit fast 40 Jahren
beschäftigt, habe ich zuweilen in einzelnen freien Stunden wieder nachgedacht. Ich
meine die ersten Gründe der Geometrie. Ich weiß nicht, ob ich Ihnen je von
meinen Ansichten darüber gesprochen habe. Auch hier habe ich manches weiter
gefestigt, und meine Überzeugung, dass wir die Geometrie nicht vollständig a priori begründen können, ist, wenn möglich, noch fester geworden.
Unterdessen werde ich wohl noch lange nicht dazu kommen, meine \textls{sehr ausgedehnten} Untersuchungen darüber für die Veröffentlichung auszuarbeiten, und vielleicht wird dies auch zu meinen Lebzeiten nie geschehen,
da ich das Geschrei der Böotier scheue, wenn ich meine Ansicht ganz aussprechen wollte. --- Seltsam ist es aber, dass es \textls{außer} der bekannten Lücke in
Euklids Geometrie, die man bisher vergeblich zu füllen versucht hat und nie
füllen wird, noch einen anderen Mangel darin gibt, den meines
Wissens bisher niemand beanstandet hat und der keineswegs leicht
(obwohl möglich) zu beheben ist. Dies ist die Definition des \textls{Planum} als einer Fläche,
in der die gerade Linie, die \textls{irgend zwei} Punkte verbindet, \textls{ganz} liegt. Diese
Definition enthält \textls{mehr}, als zur Bestimmung der Fläche nötig ist, und setzt \emph{stillschweigend} ein \textls{Theorem} voraus, das erst bewiesen werden muss. .\ .\ .\ .\ .\ .

\rule{4cm}{0.5pt}

